% Tangente, Normale, Schnittwinkel – bank/tangente-normale-schnittwinkel/ (zusammenbau v0.4, Heft)
\einheitenkopf[t4]{Tangente, Normale, Schnittwinkel}
\setcounter{aufgabe}{42}

% Anlauf (ohne Prüfkennung)
\begin{aufgabe}{Tangentengleichung – Anlauf}
\begin{teile}
\teil Gesucht ist die Tangente an den Graphen von $f$ im Punkt $P(2 | f(2))$. Was ist gegeben? \\ \kreuz{ein Punkt} \\ \kreuz{eine Steigung} \\ \kreuz{ein Winkel}
\teil $f(x) = x^3 - 2x$ – Tangente im Punkt $(2 | f(2))$? $t(x) =$ \leerfeld
\teil Liegt $Q(1 | 1)$ auf dem Graphen von $f(x) = 2x^3 - x$? Wenn ja: Tangente in $Q$? $t(x) =$ \leerfeld
\end{teile}
\end{aufgabe}

\begin{aufgabe}{Tangentengleichung – Prüfungsaufgaben}
\begin{teile}
\teil $f(x) = x^3 - 3x^2 - 2$, Tangente im Punkt $P(3 | -2)$ – Gleichung, rechnerisch? \hfill (Abitur 2025 GK) \\ $t(x) =$ \leerfeld
\teil $f(x) = 2e^x - 3$ – Gleichung der Tangente $t$ im Schnittpunkt des Graphen mit der $y$-Achse? \hfill (Abitur 2026 GK) Kontrolle: $y = 2x - 1$. $t(x) =$ \leerfeld
\teil Der Wirkstoffspiegel im Blut wird durch $k(x) = 4x \cdot e^{-0{,}1x} + 30$ beschrieben ($x$ in Tagen seit Behandlungsbeginn, $k(x)$ in Prozent des Zielwerts); es gilt $k'(x) = (4 - 0{,}4x) \cdot e^{-0{,}1x}$. Nach $15$ Tagen wird das Pflaster entfernt, danach verläuft der Spiegel entlang der Tangente $g$ im Punkt $(15 | k(15))$. Gleichung von $g$ (auf zwei Stellen) und ab welchem Tag gilt $g(x) \le 30$? \hfill (Abitur 2019 GK) \\ $g(x) \approx$ \leerfeld , ab Tag \leerfeld
\teil $f(x) = 4 - 3\mathrm{sin}(x)$ – Tangente im Punkt $(0 | f(0))$? \hfill (Abitur 2018 GK) \\ $t(x) =$ \leerfeld
\teil $f(x) = (x^2 - 5) \cdot e^x$ mit $f'(x) = (x^2 + 2x - 5) \cdot e^x$: Weise nach, dass der Graph an der Stelle $x = -1 + \sqrt{6}$ eine waagerechte Tangente hat, und gib den Funktionswert dort an (auf zwei Stellen). \hfill (Abitur 2023 GK) \\ $f \approx$ \leerfeld
\teil Der Graph von $f$ ist symmetrisch zur $y$-Achse. Die Tangente im Punkt $(2 | f(2))$ hat die Gleichung $y = -3x + 5$. Gleichung der Tangente im Punkt $(-2 | f(-2))$, mit Begründung? \hfill (Abitur 2023 GK) \\ $y =$ \leerfeld
\end{teile}
\end{aufgabe}

% Anlauf (ohne Prüfkennung)
\begin{aufgabe}{Berührung nachweisen – Anlauf}
\begin{teile}
\teil An der Stelle $x = 1$ gilt $f(1) = g(1)$, und die Gerade $g$ hat die Steigung $f'(1)$. Wie verhält sich $g$ zum Graphen von $f$ an dieser Stelle? \\ \kreuz{berührt} \\ \kreuz{schneidet nur} \\ \kreuz{weder noch}
\teil Ist $g(x) = -x + 1$ Tangente an den Graphen von $f(x) = x^3 - 2x^2 + 1$ an der Stelle $x = 1$? ja / nein: \leerfeld
\teil $f(x) = x^2 + 1$ und $g(x) = -x^2 + 4x - 1$ – weise nach, dass sich die Graphen in ihrem gemeinsamen Punkt berühren, und gib die gemeinsame Tangente an. $y =$ \leerfeld
\end{teile}
\end{aufgabe}

\begin{aufgabe}{Berührung nachweisen – Prüfungsaufgaben}
\begin{teile}
\teil Eine Rodelbahn wird für $0 \le x \le 15$ durch $h(x) = 40 - \frac{1}{10}x^2$ beschrieben ($1$ LE = $1$ m). Im Punkt $(10 | h(10))$ soll sie ohne Knick geradlinig fortgesetzt werden. Gleichung der Geraden? \hfill (Abitur 2018 GK) \\ $y =$ \leerfeld
\teil Die Tangente an den Graphen von $f(x) = x^3 - 9x^2 + 24x - 16$ im Punkt $S(0 | f(0))$ hat einen weiteren gemeinsamen Punkt mit dem Graphen. Koordinaten? \hfill (Abitur 2023 GK) \\ \leerfeld
\teil Die Abbildung zeigt den Graphen von $f$ und den Punkt $Q(0 | -1)$. Lege mit dem Lineal eine Tangente an den Graphen, die durch $Q$ geht, zeichne sie ein und gib ihre Gleichung an. \hfill (Abitur 2020 GK) \\

\begin{ksys}[xmin=-4,xmax=4,ymin=-2,ymax=6] \funktion{0.5*\x^2+1}{f} \punkt{0}{-1}{Q} \end{ksys}
$y =$ \leerfeld
\teil Der Graph von $f(x) = \frac{1}{8}x^3 - \frac{3}{4}x^2 + 3$ ist symmetrisch zu seinem Wendepunkt $W(2 | 1)$. Betrachtet werden Geraden durch $W$ mit negativer Steigung. Für welche Steigungen hat eine solche Gerade mit dem Graphen nur den Punkt $W$ gemeinsam? \hfill (Abitur 2019 GK) \\ $m$: \leerfeld
\teil Der Graph von $f(x) = -\frac{1}{3}x^3 + x^2 + \frac{4}{3}$ ist symmetrisch zu seinem Wendepunkt $W(1 | 2)$. Betrachtet werden Geraden durch $W$ mit positiver Steigung. Für welche Steigungen hat eine solche Gerade mit dem Graphen nur den Punkt $W$ gemeinsam? \hfill (Abitur 2019 GK) \\ $m$: \leerfeld
\end{teile}
\end{aufgabe}

% Anlauf (ohne Prüfkennung)
\begin{aufgabe}{Winkel – Anlauf}
\begin{teile}
\teil Gesucht ist die Stelle, an der ein Hang unter $12^\circ$ ansteigt. Was ist gegeben? \\ \kreuz{ein Punkt} \\ \kreuz{eine Steigung} \\ \kreuz{ein Winkel}
\teil Eine Gerade hat die Steigung $2$ – Steigungswinkel? $\alpha \approx$ \leerfeld °
\teil $f(x) = \frac{1}{3}x^3 - x^2$ – Steigungswinkel des Graphen im Punkt $(3 | f(3))$? $\alpha \approx$ \leerfeld °
\end{teile}
\end{aufgabe}

\begin{aufgabe}{Winkel – Prüfungsaufgaben}
\begin{teile}
\teil $f(x) = (2x + 1) \cdot e^{-x}$ mit $f'(x) = (1 - 2x) \cdot e^{-x}$ – Gleichung der Tangente $t$ im Punkt $P(1 | f(1))$ und Winkel, unter dem $t$ die $x$-Achse schneidet? \hfill (Abitur 2020 GK) \\ $t(x) \approx$ \leerfeld , Winkel: \leerfeld °
\teil $f(x) = \frac{1}{8}x^3 + 1$ – Gleichung der Tangente im Punkt $(2 | f(2))$ und Winkel, unter dem sie die $x$-Achse schneidet? \hfill (Abitur 2026 GK) \\ $t(x) =$ \leerfeld , Winkel: \leerfeld °
\teil Unter welchem Winkel schneidet der Graph von $s(x) = e^{2x}$ die Gerade $y = 2$? \hfill (Abitur 2023 GK) \\ \leerfeld °
\teil Die Graphen von $f(x) = (x + 2) \cdot e^{-\frac{1}{4}x}$ und $g(x) = 2x + 2$ schneiden sich im Punkt $S(0 | 2)$. Unter welchem Winkel schneiden sich die Tangente an den Graphen von $f$ in $S$ und die Gerade $g$? \hfill (Abitur 2018 GK) Kontrolle: $f'(x) = (\frac{1}{2} - \frac{1}{4}x) \cdot e^{-\frac{1}{4}x}$. \leerfeld °
\teil Die Graphen von $f(x) = \frac{1}{3}x^3$ und seiner Ableitung $f'$ schneiden sich an der Stelle $x = 3$. Unter welchem Winkel? \hfill (Abitur 2022 GK) \\ \leerfeld °
\teil Ein Brückenbogen einer Spielzeugbahn folgt $f(x) = -\frac{1}{40}x^2 + \frac{1}{2}x$ ($1$ LE = $1$ cm). Die Tangenten an den Stellen $2$ und $16$ dürfen einen Winkel von höchstens $30^\circ$ einschließen, sonst kippen die Wagen. Wird die Vorgabe eingehalten? \hfill (Abitur 2017 GK) \\ ja / nein: \leerfeld
\teil Die Kante eines Dachs wird für $0 \le x \le 2$ durch $f(x) = (x - 2)^2 \cdot e^{-x}$ beschrieben; es gilt $f'(x) = (-x^2 + 6x - 8) \cdot e^{-x}$. In welchem Punkt hat die Kante ein Gefälle von $30^\circ$? Auf zwei Stellen. \hfill (Abitur 2017 GK) \\ $Q$: \leerfeld
\teil Der Querschnitt eines Knaufs wird von $f(x) = 9 - x^2$ und $g(x) = -f(x)$ für $0 \le x \le 3$ begrenzt; die Spitze liegt in $S(3 | 0)$. Vorgabe: Der Winkel an der Spitze beträgt mindestens $155^\circ$. Prüfe rechnerisch. \hfill (Abitur 2024 GK) \\ \leerfeld
\teil Der Querschnitt eines Griffs wird von $f(x) = 2 - \frac{1}{8}x^2$ und $g(x) = -f(x)$ für $0 \le x \le 4$ begrenzt; die Spitze liegt in $S(4 | 0)$. Vorgabe: Der Winkel an der Spitze beträgt mindestens $100^\circ$. Prüfe rechnerisch. \hfill (Abitur 2024 GK) \\ \leerfeld
\teil $f(x) = \frac{1}{4}x^3 - x$; $k$ ist die Tangente an den Graphen im Punkt $(2 | 0)$, $g$ die Gerade $y = \frac{1}{3}x$. Steigungswinkel von $k$? Prüfe die Aussage: Für den Schnittwinkel $\alpha$ von $k$ und $g$ gilt $\mathrm{sin}\,\alpha = \mathrm{cos}\,\alpha$. \hfill (Abitur 2025 GK) \\ \leerfeld
\teil $f(x) = x^3 - 4x$; $k$ ist die Tangente an den Graphen im Punkt $(2 | 0)$, $g$ die Gerade $y = 2x$. Steigungswinkel von $k$? Prüfe die Aussage: Für den Schnittwinkel $\alpha$ von $k$ und $g$ gilt $\mathrm{sin}\,\alpha = \mathrm{cos}\,\alpha$. \hfill (Abitur 2025 GK) \\ \leerfeld
\end{teile}
\end{aufgabe}

% Anlauf (ohne Prüfkennung)
\begin{aufgabe}{Dreiecke und Figuren – Anlauf}
\begin{teile}
\teil Eine Tangente und die beiden Koordinatenachsen schließen ein Dreieck ein. Gesucht: Wie lang ist sein Rand? Was ist gefragt? \\ \kreuz{Fläche} \\ \kreuz{Umfang} \\ \kreuz{eine Eigenschaft}
\teil Die Gerade $g(x) = -3x + 6$ schließt mit den Koordinatenachsen ein Dreieck ein. Flächeninhalt? $A =$ \leerfeld
\teil $f(x) = 4 - \frac{1}{4}x^2$ – die Tangente im Punkt $(4 | 0)$ und die Koordinatenachsen begrenzen ein Dreieck. Flächeninhalt und Umfang? $A =$ \leerfeld , $U \approx$ \leerfeld
\end{teile}
\end{aufgabe}

\begin{aufgabe}{Dreiecke und Figuren – Prüfungsaufgaben}
\begin{teile}
\teil Die Kante eines Dachelements folgt $f(x) = (2 - x) \cdot e^x$ für $-2 \le x \le 0$ ($1$ LE = $5$ m). Statt einer Wand von etwa $58$ m² unter dem Graphen wird eine kleinere Wand gebaut, die von den Koordinatenachsen und der Tangente im Punkt $R(0 | 2)$ begrenzt wird. Tangentengleichung und eingesparte Fläche in m²? \hfill (Abitur 2017 GK) \\ $t(x) =$ \leerfeld , gespart: \leerfeld[m²]
\teil $f(x) = (x + 3) \cdot e^{-x}$ mit $f'(x) = -(x + 2) \cdot e^{-x}$. Die Tangente im Punkt $(0 | f(0))$ bildet mit den Koordinatenachsen ein Dreieck. Umfang, und welcher Punkt hat von allen drei Ecken denselben Abstand? \hfill (Abitur 2022 GK) \\ $U \approx$ \leerfeld , Punkt: \leerfeld
\teil $f(x) = 3 - \mathrm{sin}(x)$; die Tangenten in den Punkten $(2k\pi | f(2k\pi))$ mit $k \in \mathbb{N}$ schließen jeweils mit den Koordinatenachsen ein Dreieck ein. Begründe, dass jedes dieser Dreiecke gleichschenklig ist. \hfill (Abitur 2019 GK)
\teil $h(x) = -2x \cdot e^{-0{,}2x}$; $t$ ist die Tangente an den Graphen im Ursprung. $t$, die $x$-Achse und die Gerade $x = b$ mit $b > 0$ schließen eine Fläche von $49$ FE ein. Berechne $b$. \hfill (Abitur 2021 GK) \\ $b =$ \leerfeld
\end{teile}
\end{aufgabe}
